\documentclass[11pt]{article}

\usepackage{amsmath,amssymb,amstext,amsthm}
\usepackage{geometry}
\usepackage{fancyhdr}
\usepackage[pdftex]{graphicx}
\usepackage{xcolor}

\geometry{
  verbose,
  dvips,
  width=430pt, marginparsep=0pt, marginparwidth=0pt,
  top=90pt, headheight=12pt, headsep=20pt, footskip=30pt, bottom=80pt}

\setlength{\parskip}{\medskipamount}
\setlength{\parindent}{0pt}

\linespread{1.05}

\newtheorem{theorem}{Theorem}
\newtheorem{lemma}[theorem]{Lemma}
\newtheorem{proposition}[theorem]{Proposition}
\newtheorem{corollary}[theorem]{Corollary}
\newtheorem{conjecture}[theorem]{Conjecture}
\newtheorem{obs}{Observation}
\newtheorem{clm}{Claim}
\newtheorem{ques}{Question}
\newtheorem{problem}{Problem}

\newtheorem{ex}{Example}


\newcommand{\spc}{\hspace{0.08in}}
\newcommand{\vs}{\vspace*{.1in}}
\newcommand{\E}{\mathbb{E}}
\newcommand{\Prob}{\mathbb{P}}
\newcommand{\edit}[1]{\emph{\textcolor{blue}{#1}}}


\renewenvironment{proof}{{\noindent \textbf{\textit{Proof.}}}}
{\hfill $\Box$ \vspace*{0.1in}}
\newenvironment{proofc}{{\noindent \textit{Proof of Claim.}}}
{\hfill $\Box$}


\begin{document}

\begin{LARGE}\begin{center}\bf 12.6 Cylinders and Quadric Surfaces \end{center}
\end{LARGE}

\vs\vs



\section{Warm Up}
\textbf{Problem 1:} Graph the curve $y = 9-x^2$.

\textbf{Question 2:} What is an ellipse? What is the equation of an ellipse?



\section{Motivation}
    In 2-dimension we spend many years of math classes learning how to graph curves on the xy plane. In order to understand the tools of calculus in 3 dimensions, we similarly need to be proficient at graphing and visualizing surfaces in the xyz coordinate system.




\section{Definitions \& Theorems}
\begin{enumerate}
    \item If we just took a screenshot of a surface on any plane parallel to our coordinate planes (xy,xz,yz) we would get a picture which looks like a curve in 2-dimensions. We call this a \textbf{trace}. A trace can be thought of as a cross section.
    \item  A \textbf{cylinder} is a surface which has 2 variables and thus, its trace remains unchanged as it moves through the axis of the variable missing from the surface.
    \item A \textbf{Quadric Surface} is a surface with largest degree 2. The most general equation for this is,

    \begin{equation*}
        Ax^2+By^2+Cz^2+Dxy+Eyz+Fxz+Gx+Hy+Iz+J=0
    \end{equation*}

These will be some of the surfaces we \emph{do} want to feel comfortable drawing in 3-dimensions.

\item An \textbf{Ellipsoid} is of the form 

\begin{equation*}
    \frac{x^2}{a^2}+\frac{y^2}{b^2}+\frac{z^2}{c^2}=1.
\end{equation*} 
Its $x$-intercepts are at $\pm a$, its $y$-intercepts are at $\pm b$ and its $z$-intercepts are at $\pm c$. Here are ways to identify the ellipsoid.
\begin{itemize}
    \item ALL positive.
    \item All variables squared.
    \item Has a constant.
\end{itemize}

\item A \textbf{1-sheet Hyperboloid} is of the form 

\begin{equation*}
    \frac{x^2}{a^2}+\frac{y^2}{b^2}-\frac{z^2}{c^2}=1.
\end{equation*} 
It always is in the direction of the variable with the minus sign. Here are ways to identify the 1-sheet hyperboloid.
\begin{itemize}
    \item Has one minus sign.
    \item All variables squared.
    \item Has a constant.
\end{itemize}

To start graphing set the negative variable equal to 0 and $\pm\sqrt{\text{denominator}}$ for 3 traces (which will always be circles or ellipses).

\item A \textbf{2-sheet Hyperboloid} is of the form 

\begin{equation*}
    -\frac{x^2}{a^2}-\frac{y^2}{b^2}+\frac{z^2}{c^2}=1.
\end{equation*} 
It always is in the direction of the variable with the plus sign. Here are ways to identify the 2-sheet hyperboloid.
\begin{itemize}
    \item Has one plus sign.
    \item All variables squared.
    \item Has a constant.
\end{itemize}

To start graphing set both the negative variable equal to 0 to find the intercepts of the axis. Set the plus variable equal to some numbers divisible by the denominator.

\item A \textbf{Cone} is of the form 
\begin{equation*}
    \frac{x^2}{a^2}+\frac{y^2}{b^2}-\frac{z^2}{c^2}=0.
\end{equation*} 

It always is in the direction of the variable with the minus sign. Here are ways to identify the 1-sheet hyperboloid.
\begin{itemize}
    \item Has one minus sign.
    \item All variables squared.
    \item NO constant.
\end{itemize}

Our cone will hit the origin. To graph traces, set negative variable equal to numbers which are divisible by the denominator (traces will be circles and ellipses).

\item A \textbf{Paraboloid} is of the form

\begin{equation*}
    \frac{x^2}{a^2}+\frac{y^2}{b^2}=cz
\end{equation*}

To tell we have a paraboloid check for the following.

\begin{itemize}
    \item 3 variables 2 of which are squared
    \item squared variables are positive
\end{itemize}

Paraboloids travel in the direction of the degree one variable and opens in the direction depending on the variable $c$ (positive or negative). To graph set the degree one variable equal to 0 (unless shifted).

\item A \textbf{Hyperbolic Paraboloid} is of the form

\begin{equation*}
    \frac{x^2}{a^2}-\frac{y^2}{b^2}=cz
\end{equation*}

To tell we have a paraboloid check for the following.

\begin{itemize}
    \item 3 variables 2 of which are squared
    \item one squared variables is positive
\end{itemize}

Paraboloids travel in the direction of the degree one variable. Plug in positives and negatives into the degree one variable to check whether we get hyperbolas on certain axis.

\end{enumerate}




\section{Examples}

\begin{problem}
    Sketch the surface $z = 9-x^2$.
\end{problem}

\vspace{7cm}

\begin{problem}
    Sketch the surface $yz=1$.
\end{problem}

\vspace{5cm}


\begin{problem}
    Sketch the surface $4x^2+16y^2+z^2=64$.
\end{problem}

\vspace{7cm}

\begin{problem}
    Sketch the curve $9x^2+9z^2-4y^2=36$.
\end{problem}

\vspace{7cm}

\begin{problem}
    Sketch the curve $-z^2+9x^2+y^2+9=0$
\end{problem}

\vspace{6.5cm}

\begin{problem}
    Sketch the surface $y^2 = 4x^2+16z^2$.
\end{problem}

\vspace{6.5cm}

\begin{problem}
    Sketch the surface $z = x^2+4y^2-4$.
\end{problem}

\vspace{6.5cm}

\begin{problem}
    Sketch the surface $x = 4-5y^2-9z^2$.
\end{problem}

\vspace{9cm}

\begin{problem}
    Sketch the surface $x^2-y^2-z=0$.
\end{problem}

\newpage

\section{Scratch Work}

\end{document} 